\begin{lemma}[Missing colors telescope along PEs at one support]
Let \(u\in V(T_{q+1})\), and let
\[
 b_0<b_1<\cdots<b_p\leq q
\]
be a nonempty, strictly increasing enumeration of exactly the zero-based
rounds through \(q\) whose positive slots are PEs supported at \(u\).
Let \(\gamma_{b_0+1}=\gamma_0\) and
\(\delta_{b_p+1}=\delta_p\).  Then
\[
 \overline\varphi_{b_0}(u)
   =\bigl(\overline\varphi_{b_p+1}(u)\setminus\{\delta_p\}\bigr)
      \cup\{\gamma_0\},                                        \tag{1}
\]
and
\[
 \overline\varphi_{b_p+1}(u)
   =\bigl(\overline\varphi_{b_0}(u)\setminus\{\gamma_0\}\bigr)
      \cup\{\delta_p\}.                                       \tag{2}
\]
The enumeration, its strict increase, and the exact-range equivalence are
explicit hypotheses; \(\gamma_0\) and \(\delta_p\) are identified with
their optional data fields by explicit equalities.
\end{lemma}
\begin{proof}
For an enumerated round \(b_t\), set
\[
 A_t=\overline\varphi_{b_t}(u),\qquad
 B_t=\overline\varphi_{b_t+1}(u),
\]
and abbreviate the colors in its positive PE slot by
\(\gamma_t,\delta_t\).  We first prove the exact one-round replacement
formula
\[
 B_t=(A_t\setminus\{\gamma_t\})\cup\{\delta_t\},\qquad
 A_t=(B_t\setminus\{\delta_t\})\cup\{\gamma_t\}.               \tag{3}
\]

The exact-range hypothesis gives the PE kind and supporting-vertex
equalities at \(b_t.\mathrm{succ}\).  Eliminate the corresponding
certificate from \noderef{TashkinovSeries}; only the PE constructor of
\noderef{TashkinovSeriesStep} has that kind.  Call its maximum-defective
selected coloring \(\pi_t\) and its second stable coloring \(\rho_t\).
The selected-stability field of \noderef{MaximumDefectiveChoice} and the
second-stability field of the PE constructor, interpreted with the
vertex clause of \noderef{StableColoring}, give
\[
 \overline\pi_t(u)=A_t,
 \qquad \overline\rho_t(u)=\overline\pi_t(u).                  \tag{4}
\]
The rule \noderef{ParallelExtensionColor} says that \(\gamma_t\) is
missing at \(u\) under \(\pi_t\), hence also under \(\rho_t\).
The selected connecting edge is incident with \(u\) and has color
\(\delta_t\) under \(\pi_t\).  Since the second stability controls the
enlarged set containing \(\delta_t\), its edge clause fixes this exact
edge under \(\rho_t\).  Thus \(\delta_t\) is present at \(u\) under
\(\rho_t\), while \(\gamma_t\) is missing there; in particular the two
colors are distinct.

The PE constructor defines \(\varphi_{b_t+1}\) by
\noderef{TwoColorComponentSwap} on the
\((\gamma_t,\delta_t)\)-component containing \(u\).  Every incident
edge at \(u\) carrying either color belongs to that component.  Properness
allows at most one incident edge of either color.  The swap therefore
turns the unique \(\delta_t\)-edge into a \(\gamma_t\)-edge and changes
no incidence involving any other color.  In view of (4), this removes
\(\gamma_t\) from the missing set at \(u\), inserts \(\delta_t\), and
proves both equalities in (3).

There is no PE supported at \(u\) strictly between consecutive entries
\(b_t,b_{t+1}\), by the exact-range and strict-increase hypotheses.
Applying \noderef{HkuUnsupportedVertexColors} over that intervening
interval (or observing definitional equality when the rounds are
consecutive) yields
\[
 A_{t+1}=B_t.                                                    \tag{5}
\]
The required tree-membership premise holds because a PE supporting
vertex belongs to its current tree, and ordered extensions preserve
vertices by \noderef{TreeSequenceInitialSegment} and
\noderef{TreeSequenceVertexSet}.

By \noderef{SupportingPEColorChain},
\(\gamma_{t+1}=\delta_t\).  Define
\(R=A_0\setminus\{\gamma_0\}\).  We claim by induction on \(t\) that
\[
 A_t\setminus\{\gamma_t\}=R
 \quad\hbox{and}\quad
 B_t\setminus\{\delta_t\}=R.                                  \tag{6}
\]
For \(t=0\), this is (3).  If it holds at \(t\), then (5) and (3) give
\(A_{t+1}=R\cup\{\delta_t\}=R\cup\{\gamma_{t+1}\}\).
Moreover \(\delta_t\notin R\), because
\(R=B_t\setminus\{\delta_t\}\).  Deleting
\(\gamma_{t+1}=\delta_t\) therefore recovers exactly \(R\), and the
next instance of (3) also gives
\(B_{t+1}\setminus\{\delta_{t+1}\}=R\).  This proves (6).

At \(t=p\), (3) and (6) say
\[
 B_p=R\cup\{\delta_p\},\qquad
 A_0=R\cup\{\gamma_0\}.
\]
The endpoint colors are exactly the explicit gamma and delta fields in
the statement.  Substituting the definition of \(R\) gives (2), while
using \(R=B_p\setminus\{\delta_p\}\) gives (1).  These are the two
asserted transport identities.
\end{proof}
